% Chapter 1

\chapter{Introduction: Researching A Better Toolset} % Main chapter title

\label{Chapter1} % For referencing the chapter elsewhere, use \ref{Chapter1} 

\lhead{Chapter 1. \emph{Chapter Title Here}} % This is for the header on each page - perhaps a shortened title

%----------------------------------------------------------------------------------------

\section{Research Question}
My aptitude for computer technology combined with my artistic environment has resulted in me repairing a lot of machines. Computers need fixing for many reasons, but usually, they need fixing because the people making software have little experience with people whose primary skills are not related to computer use. The developers and their organizations are isolated from the broader population of people who use these devices. This is problematic, especially in the case of the arts. 

The arts and the humanities are tasked with producing both North America’s culture and its record of its culture. Computers are how we produce this work in 2013, and they are not simple or direct machines. A computer in 2013 is a running compromise between software that works pretty well most of the time and hardware that is generally reliable. To use these tools is a specific skill set. Using them well is as specific as using a paintbrush well. The software developed from that skillset needs to compromise. The number of people for whom computer use will be a delight in and of itself is as limited a group as for any other specialized skill. This is problematic, because art production and artistic project concepts is difficult even without the boundaries raised by software tool challenges. The more limited and specific the skill set required to use contemporary software tools, the more difficult it is to include a diversity of voices in the cultural production of genuinely contemporary work. 

When artists are excluded from technology, culture splits on lines of privilege. There are then artists, who make art, and technologists, who make technology but do not see themselves as particularly creative, or, more to the point, particularly involved in any responsibility for the ideas encoded in their work. Technology is not neutral. It is authored, and where there is authorship, there is a responsibility for ideas. When large groups are left out of communication media, particularly those tasked with producing the language with which culture speaks to itself, there comes a disconnect in our sense of self, and how we are publicly represented. 

Video games are the newest of the cultural production engines. A good game, as described by the MDA, first involves engaging mechanics – the loop systems of reward and scoring – then graphics to create a world, then sound to fill out that world. The controls are weighted this way and that, but the most profitable games - referred to as AAA or triple-A properties - are presently power fantasies. There are other fantasies out there than the ones at the end of a gun.

The value of a broad range of voices in any cultural practice should be self-evident, but video games are presently unique in their high-level uniformity. Unlike any other discipline, the basic tools for games tend to be expensive, and the toolsets that exist to generate games support specific mechanics in similar-seeming worlds. This means that art games are unfairly compared to commercial games produced with the same media. The commercial games set the ground of the conversation, which denies the subjectivity of the producer of the second-stage tools. These producers - developers and game designers - are disappeared into the system, their work only revealed by which elements an artist can or cannot subvert. This is problematic, as the second-stage producer in large part controls what the final creative director can make happen. Without cooperative tools, it is challenging to make new things.

The hardware restrictions of new games are less than advertised. While Moore’s Law is slowing down, reducing the expected increase in brute speed of linear processes in central processing units, input devices are thriving. The new research in games is no longer how to get the best graphics, but rather in how to get the most interesting experience for the investment.

The problem of interesting experiences is not trivial. Video games offer an economically advantageous distraction engine, a way to enact an artificial life. Allowing a diverse range of voices easy access to portray their own games, their own alternate or idealized modes of being, is a way of making those voices more real, of offering an alternate human experience to the “asshole simulator” {\cite Bissell} genres manufactured at much higher budgets.

My research question is to ask how to make some of those tools accessible, for people who do not want to make conventional digital games. I am basing my research on the Twine engine, popular for authoring branching text narratives, and adding an element of irony by including video, to make a video||game. This is an exploration of what a good creative tool can be, how is it different from a bad tool or an inaccessible tool.

My secondary research questions address why is it important to produce polyphony of artistic experience, and how to encourage more diverse voices in this new area of cultural production. How does working collaboratively affect tool development, and what alternate paths to learning are provided by this style of work? What are the underlying assumptions that should be included in code? What ways of working are most reproducible? Is code itself a creative practice, and how can that creative practice be described?

%----------------------------------------------------------------------------------------

\section{Initial Approach}
My initial approach and research method to this work is based on the Agile Manifesto, a software development methodology that opposes siloed, top-down software development. I have paired Agile development with the work of Helene Cixous, whose “Laugh of the Medusa” provided a template for ecriture feminine, an argument for women writing of their own experience in order to be made visible. I have also examined works such as Vera Frenkel's \textit{String Games} and some of the history of conceptual video art within Canada. 

This work is related to various texts of feminist and woman-oriented cybertheory that have appeared in the years since: Haraway’s Cyborg Manifesto, TIQQUN’s Preliminary Materials Towards A Theory of the Young-Girl. All of these are academic constructions of femininity as it is seen in relation to technology: they are feminist in the formal sense of the word. There are other senses of the term, which I will not be examining within the paper. Rather than expressing this work in context with Cixous as \textit{écriture feminine}, I will be using Cixous as a reference for the idea of the alien perspective as a perspective of resistance within a means of expression controlled by a neutral-to-hostile majority perspective. 

This approach addresses women as an alien construct to the more conventional world of technology, which has been recently associated with a masculinist performance that is unnecessary for the pure structure of good rules and the development, through that, of good software. 

My initial approach is to pair with an artist who had a game idea, take that idea, and then make it reproducible. Reproduction is, after all, the province of cyborgs: we control our biological systems, and in doing so, we have conquered what was once a hard-built destiny. By making the consequences of biological sex into a more pliant construct of gender, we have transformed what we must be to what we might be.

This work has been difficult, but I believe it to be important, not so much for the software itself - a proof of concept - but because it is important to provide software that permits people to access new technology. I am not alone in thinking this. The Arduino project, a microcontroller designed to make electronics more accessible to artists, and the Processing project, a simplified version of Java intended to improve the experience of scripting visual effects for artists, are both dedicated to the same ideals. Underlying both projects, as well as the broader Maker movement, is an ideal of participation in one’s own work.

Put simply, people who write code, particularly using the Agile methodology, are engaging in creative practice themselves, and they then display that creative practice through the artists who repurpose their work. This is a different design pattern than technology conventionally pursues, where the work is designed in isolation and released. I am using the Agile method to develop software because it does not require that one knows what the end shape of the software will be in order to pursue the end goal. Agile requires instead that developers pursue goals in sight, always keeping their development loose enough that they can repurpose their work without much effort, and it is ideal for working with artists.

What might be is a developmental model for software that permits transition in scope from the singular, minimum-viable-product model, to a model that builds on itself until complete. A small, perfect thing that does one thing very well, which permits artists to pair their own practice with a software built expressly to make their lives easier, for not too much money. This is important, because my initial approach assumes artists to be undercompensated for their work. Although artists are the central agents of production of all the cultural capital - the invisible value - of the culture industry, they are not the prime beneficiaries of the financial system that backs, stores, and distributes that capital. Therefore, software for artists needs to be inexpensive, and set up to be almost trivially easy to use. This reserves the value of scarcity to the ability of the artist, rather than applying the majority value to the role of the engineer. This kind of invisibility is the invisibility of good management, of any type of good administration. Like housekeeping, code recedes until something goes wrong.

To test this idea, I have approached people to produce video|games with the screenPerfect software in the context of a voluntary game jam – a type of collaborative space where participants work with digital tools to generate new, raw games in a limited window – and then compiling the results into an arcade machine for presentation. 

%----------------------------------------------------------------------------------------

\section{Getting Started with this Template}

If you are familiar with \LaTeX{}, then you can familiarise yourself with the contents of the Zip file and the directory structure and then place your own information into the `\texttt{Thesis.cls}' file. Section \ref{FillingFile} on page \pageref{FillingFile} tells you how to do this. Make sure you read section \ref{ThesisConventions} about thesis conventions to get the most out of this template and then get started with the `\texttt{Thesis.tex}' file straightaway.

If you are new to \LaTeX{} it is recommended that you carry on reading through the rest of the information in this document.

\subsection{About this Template}

This \LaTeX{} Thesis Template is originally based and created around a \LaTeX{} style file created by Steve R.\ Gunn from the University of Southampton (UK), department of Electronics and Computer Science. You can find his original thesis style file at his site, here:\\
\href{http://www.ecs.soton.ac.uk/~srg/softwaretools/document/templates/}{\texttt{http://www.ecs.soton.ac.uk/$\sim$srg/softwaretools/document/templates/}}

My thesis originally used the `\texttt{ecsthesis.cls}' from his list of styles. However, I knew \LaTeX{} could still format better. To get the look I wanted, I modified his style and also created a skeleton framework and folder structure to place the thesis files in.

This Thesis Template consists of that modified style, the framework and the folder structure. All the work that has gone into the preparation and groundwork means that all you have to bother about is the writing.

Before you begin using this template you should ensure that its style complies with the thesis style guidelines imposed by your institution. In most cases this template style and layout will be suitable. If it is not, it may only require a small change to bring the template in line with your institution's recommendations.

%----------------------------------------------------------------------------------------

\section{What this Template Includes}

\subsection{Folders}

This template comes as a single Zip file that expands out to many files and folders. The folder names are mostly self-explanatory:

\textbf{Appendices} -- this is the folder where you put the appendices. Each appendix should go into its own separate `\texttt{.tex}' file. A template is included in the directory.

\textbf{Chapters} -- this is the folder where you put the thesis chapters. A thesis usually has about seven chapters, though there is no hard rule on this. Each chapter should go in its own separate `\texttt{.tex}' file and they usually are split as:
\begin{itemize}
\item Chapter 1: Introduction to the thesis topic
\item Chapter 2: Background information and theory
\item Chapter 3: (Laboratory) experimental setup
\item Chapter 4: Details of experiment 1
\item Chapter 5: Details of experiment 2
\item Chapter 6: Discussion of the experimental results
\item Chapter 7: Conclusion and future directions
\end{itemize}
This chapter layout is specialised for the experimental sciences.

\textbf{Figures} -- this folder contains all figures for the thesis. These are the final images that will go into the thesis document.

\textbf{Primitives} -- this is the folder that contains scraps, particularly because one final image in the `Figures' folder may be made from many separate images and photos, these source images go here. This keeps the intermediate files separate from the final thesis figures.

\subsection{Files}

Included are also several files, most of them are plain text and you can see their contents in a text editor. Luckily, many of them are auxiliary files created by \LaTeX{} or BibTeX and which you don't need to bother about:

\textbf{Bibliography.bib} -- this is an important file that contains all the bibliographic information and references that you will be citing in the thesis for use with BibTeX. You can write it manually, but there are reference manager programs available that will create and manage it for you. Bibliographies in \LaTeX{} are a large subject and you may need to read about BibTeX before starting with this.

\textbf{Thesis.cls} -- this is an important file. It is the style file that tells \LaTeX{} how to format the thesis. You will also need to open this file in a text editor and fill in your own information (such as name, department, institution). Luckily, this is not too difficult and is explained in section \ref{FillingFile} on page \pageref{FillingFile}.

\textbf{Thesis.pdf} -- this is your beautifully typeset thesis (in the PDF file format) created by \LaTeX{}.

\textbf{Thesis.tex} -- this is an important file. This is the file that you tell \LaTeX{} to compile to produce your thesis as a PDF file. It contains the framework and constructs that tell \LaTeX{} how to layout the thesis. It is heavily commented so you can read exactly what each line of code does and why it is there. After you put your own information into the `\texttt{Thesis.cls}' file, go to this file and begin filling it in -- you have now started your thesis!

\textbf{vector.sty} -- this is a \LaTeX{} package, it tells \LaTeX{} how to typeset mathematical vectors. Using this package is very easy and you can read the documentation on the site (you just need to look at the `\texttt{vector.pdf}' file):\\
\href{http://www.ctan.org/tex-archive/macros/latex/contrib/vector/}{\texttt{http://www.ctan.org/tex-archive/macros/latex/contrib/vector/}}

\textbf{lstpatch.sty} -- this is a \LaTeX{} package required by this LaTeX template and is included as not all \TeX{} distributions have it installed by default. You do not need to modify this file.

Files that are \emph{not} included, but are created by \LaTeX{} as auxiliary files include:

\textbf{Thesis.aux} -- this is an auxiliary file generated by \LaTeX{}, if it is deleted \LaTeX{} simply regenerates it when you run the main `\texttt{.tex}' file.

\textbf{Thesis.bbl} -- this is an auxiliary file generated by BibTeX, if it is deleted, BibTeX simply regenerates it when you run the main tex file. Whereas the `\texttt{.bib}' file contains all the references you have, this `\texttt{.bbl}' file contains the references you have actually cited in the thesis and is used to build the bibliography section of the thesis.

\textbf{Thesis.blg} -- this is an auxiliary file generated by BibTeX, if it is deleted BibTeX simply regenerates it when you run the main `\texttt{.tex}' file.

\textbf{Thesis.lof} -- this is an auxiliary file generated by \LaTeX{}, if it is deleted \LaTeX{} simply regenerates it when you run the main `\texttt{.tex}' file. It tells \LaTeX{} how to build the `List of Figures' section.

\textbf{Thesis.log} -- this is an auxiliary file generated by \LaTeX{}, if it is deleted \LaTeX{} simply regenerates it when you run the main `\texttt{.tex}' file. It contains messages from \LaTeX{}, if you receive errors and warnings from \LaTeX{}, they will be in this `\texttt{.log}' file.

\textbf{Thesis.lot} -- this is an auxiliary file generated by \LaTeX{}, if it is deleted \LaTeX{} simply regenerates it when you run the main `\texttt{.tex}' file. It tells \LaTeX{} how to build the `List of Tables' section.

\textbf{Thesis.out} -- this is an auxiliary file generated by \LaTeX{}, if it is deleted \LaTeX{} simply regenerates it when you run the main `\texttt{.tex}' file.


So from this long list, only the files with the `\texttt{.sty}', `\texttt{.bib}', `\texttt{.cls}' and `\texttt{.tex}' extensions are the most important ones. The other auxiliary files can be ignored or deleted as \LaTeX{} and BibTeX will regenerate them.

%----------------------------------------------------------------------------------------

\section{Filling in the `\texttt{Thesis.cls}' File}\label{FillingFile}

You will need to personalise the thesis template and make it your own by filling in your own information. This is done by editing the `\texttt{Thesis.cls}' file in a text editor.

Open the file and scroll down, past all the `$\backslash$\texttt{newcommand}\ldots' items until you see the entries for `\texttt{University Name}', `\texttt{Department Name}', etc\ldots.

Fill out the information about your group and institution and ensure you keep to block capitals where it asks you to. You can also insert web links, if you do, make sure you use the full URL, including the `\texttt{http://}' for this.

The last item you should need to fill in is the Faculty Name (in block capitals). When you have done this, save the file and recompile `\texttt{Thesis.tex}'. All the information you filled in should now be in the PDF, complete with web links. You can now begin your thesis proper!

%----------------------------------------------------------------------------------------

\section{The `\texttt{Thesis.tex}' File Explained}

The \texttt{Thesis.tex} file contains the structure of the thesis. There are plenty of written comments that explain what pages, sections and formatting the \LaTeX{} code is creating. Initially there seems to be a lot of \LaTeX{} code, but this is all formatting, and it has all been taken care of so you don't have to do it.

Begin by checking that your information on the title page is correct. For the thesis declaration, your institution may insist on something different than the text given. If this is the case, just replace what you see with what is required.

Then comes a page which contains a funny quote. You can put your own, or quote your favourite scientist, author, person, etc\ldots Make sure to put the name of the person who you took the quote from.

Next comes the acknowledgements. On this page, write about all the people who you wish to thank (not forgetting parents, partners and your advisor/supervisor).

The contents pages, list of figures and tables are all taken care of for you and do not need to be manually created or edited. The next set of pages are optional and can be deleted since they are for a more technical thesis: insert a list of abbreviations you have used in the thesis, then a list of the physical constants and numbers you refer to and finally, a list of mathematical symbols used in any formulae. Making the effort to fill these tables means the reader has a one-stop place to refer to instead of searching the internet and references to try and find out what you meant by certain abbreviations or symbols.

The list of symbols is split into the Roman and Greek alphabets. Whereas the abbreviations and symbols ought to be listed in alphabetical order (and this is \emph{not} done automatically for you) the list of physical constants should be grouped into similar themes.

The next page contains a one line dedication. Who will you dedicate your thesis to?

Finally, there is the section where the chapters are included. Uncomment the lines (delete the `\texttt{\%}' character) as you write the chapters. Each chapter should be written in its own file and put into the `Chapters' folder and named `\texttt{Chapter1}', `\texttt{Chapter2}, etc\ldots Similarly for the appendices, uncomment the lines as you need them. Each appendix should go into its own file and placed in the `Appendices' folder.

After the preamble, chapters and appendices finally comes the bibliography. The bibliography style (called `\texttt{unsrtnat}') is used for the bibliography and is a fully featured style that will even include links to where the referenced paper can be found online. Do not under estimate how grateful you reader will be to find that a reference to a paper is just a click away. Of course, this relies on you putting the URL information into the BibTeX file in the first place.

%----------------------------------------------------------------------------------------

\section{Thesis Features and Conventions}\label{ThesisConventions}

To get the best out of this template, there are a few conventions that you may want to follow.

One of the most important (and most difficult) things to keep track of in such a long document as a thesis is consistency. Using certain conventions and ways of doing things (such as using a Todo list) makes the job easier. Of course, all of these are optional and you can adopt your own method.

\subsection{Printing Format}

This thesis template is designed for single sided printing as most theses are printed and bound this way. This means that the left margin is always wider than the right (for binding). Four out of five people will now judge the margins by eye and think, ``I never 
noticed that before.''.

The headers for the pages contain the page number on the right side (so it is easy to flick through to the page you want) and the chapter name on the left side.

The text is set to 11 point and a line spacing of 1.3. Generally, it is much more readable to have a smaller text size and wider gap between the lines than it is to have a larger text size and smaller gap. Again, you can tune the text size and spacing should you want or need to. The text size can be set in the options for the `$\backslash$\texttt{documentclass}' command at the top of the `\texttt{Thesis.tex}' file and the spacing can be changed by setting a different value in the `$\backslash$\texttt{setstretch}' commands (scattered throughout the `\texttt{Thesis.tex}' file).

\subsection{Using US Letter Paper}

The paper size used in the template is A4, which is a common -- if not standard -- size in Europe. If you are using this thesis template elsewhere and particularly in the United States, then you may have to change the A4 paper size to the US Letter size. Unfortunately, this is not as simple as replacing instances of `\texttt{a4paper}' with `\texttt{letterpaper}'.

This is because the final PDF file is created directly from the \LaTeX{} source using a program called `\texttt{pdfTeX}' and in certain conditions, paper size commands are ignored and all documents are created with the paper size set to the size stated in the configuration file for pdfTeX (called `\texttt{pdftex.cfg}').

What needs to be done is to change the paper size in the configuration file for \texttt{pdfTeX} to reflect the letter size. There is an excellent tutorial on how to do this here: \\
\href{http://www.physics.wm.edu/~norman/latexhints/pdf_papersize.html}{\texttt{http://www.physics.wm.edu/$\sim$norman/latexhints/pdf\_papersize.html}}

It may be sufficient just to replace the dimensions of the A4 paper size with the US Letter size in the \texttt{pdftex.cfg} file. Due to the differences in the paper size, the resulting margins may be different to what you like or require (as it is common for Institutions to dictate certain margin sizes). If this is the case, then the margin sizes can be tweaked by opening up the \texttt{Thesis.cls} file and searching for the line beginning with, `$\backslash$\texttt{setmarginsrb}' (not very far down from the top), there you will see the margins specified. Simply change those values to what you need (or what looks good) and save. Now your document should be set up for US Letter paper size with suitable margins.

\subsection{References}

The `\texttt{natbib}' package is used to format the bibliography and inserts references such as this one \citep{Reference3}. The options used in the `\texttt{Thesis.tex}' file mean that the references are listed in numerical order as they appear in the text. Multiple references are rearranged in numerical order (e.g. \citep{Reference2, Reference1}) and multiple, sequential references become reformatted to a reference range (e.g. \citep{Reference2, Reference1, Reference3}). This is done automatically for you. To see how you use references, have a look at the `\texttt{Chapter1.tex}' source file. Many reference managers allow you to simply drag the reference into the document as you type.

Scientific references should come \emph{before} the punctuation mark if there is one (such as a comma or period). The same goes for footnotes\footnote{Such as this footnote, here down at the bottom of the page.}. You can change this but the most important thing is to keep the convention consistent throughout the thesis. Footnotes themselves should be full, descriptive sentences (beginning with a capital letter and ending with a full stop).

To see how \LaTeX{} typesets the bibliography, have a look at the very end of this document (or just click on the reference number links).

\subsection{Figures}

There will hopefully be many figures in your thesis (that should be placed in the `Figures' folder). The way to insert figures into your thesis is to use a code template like this:
\begin{verbatim}
\begin{figure}[htbp]
  \centering
    \includegraphics{Figures/Electron.pdf}
    \rule{35em}{0.5pt}
  \caption[An Electron]{An electron (artist's impression).}
  \label{fig:Electron}
\end{figure}
\end{verbatim}
Also look in the source file. Putting this code into the source file produces the picture of the electron that you can see in the figure below.

\begin{figure}[htbp]
	\centering
		\includegraphics{Figures/Electron.pdf}
		\rule{35em}{0.5pt}
	\caption[An Electron]{An electron (artist's impression).}
	\label{fig:Electron}
\end{figure}

Sometimes figures don't always appear where you write them in the source. The placement depends on how much space there is on the page for the figure. Sometimes there is not enough room to fit a figure directly where it should go (in relation to the text) and so \LaTeX{} puts it at the top of the next page. Positioning figures is the job of \LaTeX{} and so you should only worry about making them look good!

Figures usually should have labels just in case you need to refer to them (such as in Figure \ref{fig:Electron}). The `$\backslash$\texttt{caption}' command contains two parts, the first part, inside the square brackets is the title that will appear in the `List of Figures', and so should be short. The second part in the curly brackets should contain the longer and more descriptive caption text.

The `$\backslash$\texttt{rule}' command is optional and simply puts an aesthetic horizontal line below the image. If you do this for one image, do it for all of them.

The \LaTeX{} Thesis Template is able to use figures that are either in the PDF or JPEG file format.

\subsection{Typesetting mathematics}

If your thesis is going to contain heavy mathematical content, be sure that \LaTeX{} will make it look beautiful, even though it won't be able to solve the equations for you.

The ``Not So Short Introduction to \LaTeX{}'' (available \href{http://www.ctan.org/tex-archive/info/lshort/english/lshort.pdf}{here}) should tell you everything you need to know for most cases of typesetting mathematics. If you need more information, a much more thorough mathematical guide is available from the AMS called, ``A Short Math Guide to \LaTeX{}'' and can be downloaded from:\\
\href{ftp://ftp.ams.org/pub/tex/doc/amsmath/short-math-guide.pdf}{\texttt{ftp://ftp.ams.org/pub/tex/doc/amsmath/short-math-guide.pdf}}

There are many different \LaTeX{} symbols to remember, luckily you can find the most common symbols \href{http://www.sunilpatel.co.uk/latexsymbols.html}{here}. You can use the web page as a quick reference or crib sheet and because the symbols are grouped and rendered as high quality images (each with a downloadable PDF), finding the symbol you need is quick and easy.

You can write an equation, which is automatically given an equation number by \LaTeX{} like this:
\begin{verbatim}
\begin{equation}
E = mc^{2}
  \label{eqn:Einstein}
\end{equation}
\end{verbatim}

This will produce Einstein's famous energy-matter equivalence equation:
\begin{equation}
E = mc^{2}
\label{eqn:Einstein}
\end{equation}

All equations you write (which are not in the middle of paragraph text) are automatically given equation numbers by \LaTeX{}. If you don't want a particular equation numbered, just put the command, `$\backslash$\texttt{nonumber}' immediately after the equation.

%----------------------------------------------------------------------------------------

\section{Sectioning and Subsectioning}

You should break your thesis up into nice, bite-sized sections and subsections. \LaTeX{} automatically builds a table of Contents by looking at all the `$\backslash$\texttt{chapter}$\{\}$', `$\backslash$\texttt{section}$\{\}$' and `$\backslash$\texttt{subsection}$\{\}$' commands you write in the source.

The table of Contents should only list the sections to three (3) levels. A `$\backslash$\texttt{chapter}$\{\}$' is level one (1). A `$\backslash$\texttt{section}$\{\}$' is level two (2) and so a `$\backslash$\texttt{subsection}$\{\}$' is level three (3). In your thesis it is likely that you will even use a `$\backslash$\texttt{subsubsection}$\{\}$', which is level four (4). Adding all these will create an unnecessarily cluttered table of Contents and so you should use the `$\backslash$\texttt{subsubsection$^{*}\{\}$}' command instead (note the asterisk). The asterisk ($^{*}$) tells \LaTeX{} to omit listing the subsubsection in the Contents, keeping it clean and tidy.

%----------------------------------------------------------------------------------------

\section{In Closing}

You have reached the end of this mini-guide. You can now rename or overwrite this pdf file and begin writing your own `\texttt{Chapter1.tex}' and the rest of your thesis. The easy work of setting up the structure and framework has been taken care of for you. It's now your job to fill it out!

Good luck and have lots of fun!

\begin{flushright}
Guide written by ---\\
Sunil Patel: \href{http://www.sunilpatel.co.uk}{www.sunilpatel.co.uk}
\end{flushright}
